\section{Operações com Matrizes}

\begin{def:adicao de matrizes}
Sejam $A$ e $B$	matrizes de tipo $m \times n$.
A soma delas será uma matriz $C$ de tipo $m \times n$ tal que, se $a_{ij}$ é entrada de
$A$ e $b_{ij}$ entrada de $B$, então a entrada $c_{ij}$ de $C$ é
\[
	c_{ij} = a_{ij} + b_{ij}
\]
\end{def:adicao de matrizes}

\begin{exp:soma matrizes}
\[
	\begin{bmatrix}
		2 & 3 \\
		4 & 5
	\end{bmatrix} +
	\begin{bmatrix}
		4 & 1  \\
		7 & 10
	\end{bmatrix} =
	\begin{bmatrix}
		2 + 4 & 3 + 1  \\
		4 + 7 & 5 + 10
	\end{bmatrix} =
	\begin{bmatrix}
		6  & 4  \\
		11 & 15
	\end{bmatrix}
\]
\end{exp:soma matrizes}

\begin{def:produto escalar matriz}
Seja $A$ uma matriz do tipo $m \times n$ e $c \in \real$.
O produto de $A$ por $c$ é uma matriz $P$ cuja entrada $p_{ij}$ é dada por
\[
	p_{ij} = c \cdot a_{ij}
\]
\end{def:produto escalar matriz}

\begin{exp:produto escalar matriz}
\[
	2 \cdot \begin{bmatrix}
		3 & 5  \\
		8 & 10
	\end{bmatrix} =
	\begin{bmatrix}
		2 \cdot 3 & 2 \cdot 5  \\
		2 \cdot 8 & 2 \cdot 10
	\end{bmatrix} =
	\begin{bmatrix}
		6  & 10 \\
		16 & 20
	\end{bmatrix}
\]
\end{exp:produto escalar matriz}

\begin{prop:propriedades soma matrizes}
A soma de duas matrizes é dotada das propriedades de
\begin{enumerate}
	\item Associatividade: $(A + B) + C = A + (B + C)$.
	\item Comutativadade: $A + B = B + A$.
	\item Elemento Neutro: $A + O = A$, onde $O$ é matriz nula.
	\item Inverso Aditivo: $A + (-A) = O$
\end{enumerate}
\end{prop:propriedades soma matrizes}
\begin{proof}
	Segue diretamente das propriedades de soma de números reais.
\end{proof}

\begin{def:produto de matrizes}
O produto $C$ de duas matrizes $A$ e $B$ de tipos $m \times n$ e $p \times q$ está definido
se $n = p$.
A entrada $c_{ij}$ de $C$ é dada por
\[
	c_{ij} = \sum_{k=1}^n a_{ik} \cdot b_{kj}
\]
\end{def:produto de matrizes}

\begin{prop:elemento neutro produto}
A matriz identidade é o elemento neutro da multiplicação de matrizes.
\end{prop:elemento neutro produto}
\begin{proof}
	Seja $A$ uma matriz $m \times n$.
	Com efeito, a entrada $c_{ij}$ de $A \cdot \mathbf I_n$ é definida como
	\[
		c_{ij} = \sum_{k=1}^n a_{ik} \lambda_{kj}
	\]
	Temos que $\lambda_{kj} = 1$ quando $k = j$, e zero caso contrário.
	Consequentemente,
	\[
		c_{ij} = a_{ij}
	\]
	para quaisquer índices $i \in [m]$ e $j \in [n]$.
	O mesmo vale analogamente para $\mathbf I_n \cdot A$.
\end{proof}

\begin{def:matriz do sistema}
Seja $S$ um sistema linear de $m$ equações em $n$ variáveis.
A forma matricial de $S$ é composta por matrizes $A \in \real^{m \times n}$,
$\vec x \in \real^n$ e $\vec b \in \real^m$ tal que
\begin{enumerate}
	\item A entrada $a_{ij}$ de $A$ é o $j$-ésimo coeficiente da $i$-ésima equação
	\item $\vec x$
\end{enumerate}
\end{def:matriz do sistema}
